\documentclass[11pt]{article}

\usepackage{amsmath,amssymb,graphicx}
\usepackage[margin=1in]{geometry}
\usepackage{booktabs}

\title{Emergent Phase Synchronization Dynamics Probe on a Frozen Kernel-Induced Cochain Operator Architecture}
\author{Tomislav Rupi\'c \\ Independent Researcher}
\date{March 2026}

\begin{document}

\maketitle

\begin{abstract}
Stages 10--16 established a frozen pre-geometric ladder on a kernel-induced cochain operator architecture: single-packet morphology, collective interaction, coarse envelope organization, scale-coupled fragility, negative weak binding tests, localized relational ordering, and a negative relational-feedback probe. Stage 17 changes the tested organizing variable. Instead of asking whether weak binding or weak relational feedback can stabilize morphology, it asks whether weak phase-coupled feedback can produce self-sustained collective timing structure before any geometric or binding law appears. A narrow nine-run pilot is executed on three representative collective regimes: a phase-ordered symmetric triad, a counter-propagating corridor, and a clustered composite anchor. Each case is run under null control, very weak phase coupling with $\varepsilon=0.01$, and weak phase coupling with $\varepsilon=0.02$. The phase-coupling signal is derived from internal packet phases on the frozen architecture and is scored using a dual-proxy anti-false-positive scheme. The result is cleanly negative. All nine runs remain in the class ``no synchronization effect'', no sustained Kuramoto-style order gain appears, no phase-variance compression is detected, no phase-cluster stabilization occurs, and no morphology correlate is observed. The bounded conclusion is that, at the tested strengths, the frozen architecture still does not exhibit an emergent collective timing channel.
\end{abstract}

\section{Introduction}

The program through Stage 16 is cumulative and deliberately restrictive. Stage 10 introduced a reproducible morphology instrument for wave-packet transport. Stage 11 extended that instrument to collective interaction structure. Stage 12 coarse-grained multi-packet fields into envelope and basin summaries. Stage 13 coupled refinement directly to interaction and coarse scales and showed that coarse organization sharpens under refinement while composite interaction labels remain fragile. Stage 14 tested several weak structural branches and found no effective binding channel at the tested strengths. Stage 15 then asked whether frozen collective data support a stable relational ordering that behaves metric-like under refinement. Localized ordering appeared, but it was regime-dependent and not anchored to the strongest composite regime. Stage 16 finally tested whether that relational ordering could act back on the dynamics as a weak organizing influence, and the answer was negative.

Stage 17 therefore changes axis rather than strength. It asks whether collective timing can appear even when weak binding and weak relational feedback do not. The question is narrow: can weak phase-coupled feedback, derived from internal phase structure on the frozen architecture, produce self-sustained synchronization regimes that stabilize collective morphology? The stage does not assume an oscillator model or a geometric interpretation. It only tests whether a phase-ordering channel exists at all.

\section{Frozen Context and Pilot Design}

All kernels, projectors, operator constructions, and packet families remain frozen. The projected transverse sector and periodic boundary condition are retained, together with the packet family inherited from the strongest coherent transverse seed: bandwidth $w=0.1$, cosine carrier, central wave number $k=1.0$, and amplitude scale $0.5$.

The Stage 17 pilot uses three representatives:
\begin{itemize}
\item phase-ordered symmetric triad,
\item counter-propagating corridor,
\item clustered composite anchor.
\end{itemize}

These span a weakly structured phase-ordered configuration, a transport-dominated control, and the strongest available composite anchor. Each case is run under three conditions:
\begin{itemize}
\item null control,
\item very weak phase coupling with $\varepsilon=0.01$,
\item weak phase coupling with $\varepsilon=0.02$.
\end{itemize}

This gives a total of nine base runs. As in earlier stages, a $(12 \rightarrow 24)$ refinement follow-up is allowed only if a base-resolution feedback run shows a real positive signal that survives anti-false-positive checks. No run satisfied that gate.

\section{Phase Proxy and Coupling Construction}

Two phase observables are used.

\subsection{Canonical proxy}

The canonical phase observable is the analytic-envelope phase. For each packet, a local scalar signal is built by weighted spatial averaging around the evolving packet centroid, and a Hilbert transform is applied to form an analytic signal. The phase is then read as the argument of that analytic signal. This proxy is local, continuous, and directly tied to packet evolution.

\subsection{Control proxy}

The control observable is the dominant-mode projection phase. The local packet state is projected onto the frozen narrowband carrier mode, and the complex argument of that projection is recorded. This proxy is more rigid and spectral, and it serves as a control against extraction-method artifacts.

The actual feedback step uses the dominant-mode phase because it is instantaneous and local. The analytic-envelope phase remains the canonical scoring signal.

The weak feedback term is of the form
\[
\dot X \;\mapsto\; \dot X + \varepsilon \sum_j W_{ij}\sin(\phi_j - \phi_i),
\]
where the weights $W_{ij}$ are derived from the frozen packet coupling graph and row-normalized locally. No global normalization is introduced, and no coefficient is tuned beyond the pre-registered weak regime.

\section{Anti-False-Positive Scoring}

Stage 17 is the first layer where false synchronization signals become easy to overread. The pilot therefore uses explicit anti-false-positive gates.

A candidate positive must satisfy all of the following:
\begin{itemize}
\item dual-proxy consistency: the canonical proxy must show a gain, and the control proxy must not contradict it strongly,
\item sustained order: gains in mean or median order parameter, or in the fraction of time above the fixed threshold $R>0.7$, rather than one-frame peaks,
\item phase-variance compression over a nontrivial time fraction,
\item stable cluster structure or longer dominant-partition dwell than null,
\item frequency-spread compression,
\item at least one morphology correlate, such as composite lifetime, binding persistence, or coarse persistence.
\end{itemize}

For the symmetric triad and the counter-propagating corridor, the bar is deliberately higher because trivial symmetry can inflate coherence measures. High order parameter alone does not count as synchronization in these cases.

The allowed output labels are:
\begin{itemize}
\item no synchronization effect,
\item transient synchronization,
\item scale-fragile phase locking,
\item scale-stable phase locking.
\end{itemize}

\section{Pilot Outcome}

The Stage 17 outcome is uniform:
\begin{itemize}
\item no synchronization effect: 9,
\item transient synchronization: 0,
\item scale-fragile phase locking: 0,
\item scale-stable phase locking: 0.
\end{itemize}

No case is promoted to the refinement ladder.

The reported deltas remain at numerical zero or near-zero throughout the pilot. There is no sustained increase in the canonical order parameter, no increase in the fraction of time above the threshold $R>0.7$, no drop in mean phase variance, no stable change in cluster dwell fraction, and no morphology correlate in any of the nine runs.

\subsection{Phase-ordered symmetric triad}

The phase-ordered triad is the most natural candidate for a synchronization response. Nevertheless, both tested feedback strengths remain indistinguishable from null at the recorded precision. The interaction label stays in the \emph{transient binding regime}, the coarse label remains \emph{diffuse coarse field regime}, and all synchronization deltas are numerically negligible.

\subsection{Counter-propagating corridor}

The corridor case is the strongest transport-style control. Again, both tested couplings remain null. No sustained order gain is observed, no phase-variance compression appears, and no morphology quantity changes relative to null.

\subsection{Clustered composite anchor}

The clustered composite anchor is retained because it had the strongest collective composite morphology in earlier stages. If a weak phase channel were going to interact with already structured collective behavior, this would be a natural place for the first hint. None appears. The case remains unchanged under both tested couplings, with no synchronization gain and no morphology response.

\section{Interpretation of the Negative Result}

The Stage 17 null is not redundant with earlier null results. It closes a different possible route to collective law. By this point the architecture has shown three things clearly:
\begin{itemize}
\item structured collective morphology can exist,
\item localized relational ordering can exist,
\item neither weak binding feedback, weak relational feedback, nor weak phase synchronization promote those structures into a new collective dynamical law.
\end{itemize}

This strengthens the emerging constraint map. The frozen architecture is rich enough to support organized morphology and local ordering, but at the tested weak couplings it remains conservative: it does not readily convert such structure into self-sustained collective law.

Stage 17 therefore adds a specific negative statement: collective timing does not emerge automatically even when the morphology already exhibits organization and even when a local phase proxy can be defined cleanly.

\section{Interpretation Boundary}

Stage 17 does not support:
\begin{itemize}
\item effective binding claims,
\item emergent spacetime claims,
\item physical oscillator claims,
\item causal-law claims.
\end{itemize}

The only bounded claims supported here are:
\begin{itemize}
\item the frozen architecture admits well-defined local phase proxies,
\item at $\varepsilon=0.01$ and $0.02$ weak phase coupling does not generate measurable collective timing structure,
\item no morphology correlate appears alongside the null synchronization read,
\item the architecture still lacks an emergent organizing law at this layer.
\end{itemize}

\section{Conclusion}

Stage 17 is the first direct probe of whether weak phase synchronization can emerge as a collective organizing axis on the frozen architecture. The answer, in the first executable pilot, is negative. Across nine base runs spanning a phase-ordered triad, a counter-propagating corridor, and a clustered composite anchor, no self-sustained timing structure appears at the tested strengths.

This result is useful because it closes another natural route to emergent law. The current architecture can sculpt form, support localized ordering, and maintain structured collective morphology, but it still does not convert those ingredients into a weak collective timing regime. The next sensible step is therefore not to increase the coupling coefficient blindly, but to test a qualitatively different coupling class, such as delayed or history-dependent influence.

\end{document}
